\chapter{Preliminaries}\label{chapt:preliminaries}
\minitoc
\lipsum[1]

\section{Number Systems}
\lipsum[1]
\subsection{Base/Radix Concepts}
\lipsum[1]
\subsection{Positional Notation}
\lipsum[1]
\subsection{Binary Arithmetic}
\lipsum[1]

\section{Number Representations}
\lipsum[1]
\subsection{Unsigned Numbers}
\lipsum[1]
\subsection{Signed Numbers}
\lipsum[1]
\subsubsection{One's Complement}
\lipsum[1]
\subsubsection{Two's Complement}
\lipsum[1]
\subsection{Overflow Detection}
\lipsum[1]

\section{Modular Arithmetic}
\lipsum[1]
\subsection{Modulo Operation}
\lipsum[1]
\subsection{Residue Classes}
\lipsum[1]

\section{Binary Coded Decimal (BCD) and Other Encodings}
\lipsum[1]
\subsection{BCD}
\lipsum[1]
\subsection{Gray Code}
\lipsum[1]

\section{Floating-Point Arithmetic}
\lipsum[1]
\subsection{Scientific Notation and Exponent Representation}
\lipsum[1]
\subsection{IEEE Floating-Point Standard}
\lipsum[1]

\section{Set Notation}
\lipsum[1]

\section{Digital Domain}

\subsection{Signals in Digital Domain}

In the electronic digital domain we work with two values only (see Figure \ref{fig1}):

\placedrawing[H]{figures/01_01_digital_0_and_1.lp}{The two levels of the signal in the digital domain. Low level (0 Volt) for {\tt 0} or {\tt false}, and high level ($V_{DD}$ Volt) for {\tt 1} or {\tt true}.}{fig1}

\begin{itemize}
\item 0, represented by the electrical value 0 V, having two meanings:
\begin{itemize}
\item the numerical value {\tt 0}
\item the logic value {\tt false}
\end{itemize}
\item 1, represented by the electrical value $V_{DD}$ V, having two meanings:
\begin{itemize}
\item the numerical value {\tt 1}
\item the logic value {\tt true}
\end{itemize}
\end{itemize}

The signals used in a digital system are of two types:
\begin{itemize}
\item non-periodic, with a certain structure ("random") adapted to the process of simulating the operation of a digital system
\item periods, usually with a strictly alternating structure, used as clock signals by which the operation of a digital system is synchronized.
\end{itemize}

Any simulation is based on the generation of waveforms with which the simulated circuit is stimulated on its inputs.
In the following System Verilog module, two waveforms are generated, one "random" and one usable as a periodic signal.

{\small
\begin{shaded*}
\begin{lstlisting}[language=Verilog, caption= , morekeywords={always_ff, always_comb, logic, loop, repeat, doInParallel, for}]
/*************************************************************************
File name: waveFormGenerator.sv
Circuit name: no circuit, only wave formes
Description: two waveforms are generated, a "random" one and a periodical
             one: a clock signal
*************************************************************************/
module waveFormGenerator();
    logic randomWave;
    logic clock     ;
    initial begin       randomWave = 0  ;
                    #2  randomWave = 1  ;
                    #6  randomWave = 0  ;
                    #4  randomWave = 1  ;
                    #8  randomWave = 0  ;
                    #5  $stop           ;
            end
    initial begin               clock = 0       ;
                    forever #2  clock = ~clock  ;
            end
endmodule
\end{lstlisting}
\end{shaded*}
}

By simulating the {\tt waveFormGenerator.sv} module, the behavior represented in Figure \ref{wave} results.



\placedrawing[H]{figures/01_02_waveSim.lp}{Waveforms generated for a Vivado System Verilog simulation}{wave}


In a digital system we can distinguish two extreme types of circuits:
\begin{itemize}
\item purely logical circuits that operate on some binary signals with some symbolic meaning (example: the transcoder for displaying the numbers from 0 to 9 using 7 segments, represented in Figure \ref{fig2}a)

\item purely numerical circuits that operate on binary signals with numerical significance (example: a modulo 16 adder for two numbers represented by 4 bits each, represented in Figure \ref{fig2}b)
\end{itemize}
Obviously, an intermediate version is always possible (example: a circuit that compares two numbers, A and B, and indicates through three logical signals of one bit each: $A>B, \; A=B, \; A\leq B$)

\placedrawing[H]{figures/01_03_7seg_and_adder.lp}{The version of digital circuits. {\bf a.} Logic circuit: trans-coder for seven-segment display. {\bf b.} Numeric circuit: four-bit numbers adder.}{fig2}

\subsection{Character encoding}

There is American Standard Code for Information Interchange (ASCII),  a character encoding standard used to represent text in computers.
ASCII codes are represented on 8-bit binary configuration: $B = \{b_7, b_6, \ldots, b_0\}$. For example: from $B = 8'b011\_0000$ to $B = 8'b011\_1001$ the codes are used to represent digital numbers from 0 to 9, while from $B = 8'b100\_0000$ to $B = 8'b011\_1111$ the codes are used to represent the letters $@, A, B, \ldots, N, O$.

See also UTF-16 (16-bit Unicode Transformation Format) which is a character encoding capable of encoding all characters  used in our computer.

\subsection{Number encoding}

While a decimal number (a number represented in base 10) represented on $n$ digits has the form:
$$D = 10^{n-1} \times d_{n-1} + 10^{n-2} \times d_{n-2} + \ldots + 10^{1} \times d_{1} + 10^{0} \times d_{0}$$
where: $d_i \in \{0, 1, \ldots, 9\}$ for $i = 0, 1, \ldots, n-1$
a binary number (a number represented in base 2) represented on $n$ bits has the form:
$$B = 2^{n-1} \times b_{n-1} + 2^{n-2} \times b_{n-2} + \ldots + 2^{1} \times b_{1} + 2^{0} \times b_{0}$$
where: $b_i \in \{0, 1\}$ for $i = 0, 1, \ldots, n-1$

Examples:
\begin{itemize}
    \item $253 = 10^2 \times 2 + 10^1 \times 5 + 10^0 \times 3$
    \item $1010 = 2^3 \times 1 + 2^2 \times 0 + 2^1 \times 1 + 2^0 \times 0 = 10$
\end{itemize}

\subsubsection{Negative numbers}

There are two main mode of representing signed numbers:
\begin{enumerate}
    \item sign \& magnitude: $signedNumber[n-1:0] = \{sgn[0], value[n-2:0] \}$
    \item 2's complement: $- signedNumber = \sim(signedNumber) + 1$
\end{enumerate}
where: $\sim$ is the operator bitwise complement (or 1's complement).

Examples:
\begin{enumerate}
    \item using a 5-bit representation (the most significant bit for sign, and the rest of bits a magnitude) $+6 = 0\_0110$, while $-6 = 1\_0110$
    \item using a 5-bit compact representation $+6 = 00110$ while $-6 = \sim(00110) + 1 = 11001 + 1 = 11010$. And if we convert back -6 to +6, then: $\sim(11010) + 1 = 00101 + 1 = 00110 = +6$
\end{enumerate}

\subsubsection{Adding binary numbers}

Adding two 1-bit numbers provides sum and carry according to Table \ref{add}.

{\small
\begin{table}[H]
  \begin{center}
    \caption{Elementary adder}
    \label{add}
    \begin{tabular}{|c|c||c|c|}
      \hline
      A&    B &  sum & carry     \\ \hline \hline
      0&  0&   0&       0    \\ \hline
      0&  1&   1&       0    \\ \hline
      1&  0&   1&       0    \\ \hline
      1&  1&   0&       1    \\ \hline
    \end{tabular}
  \end{center}
\end{table}
}
The carry signal is used in the addition of many-bit operation. It is added to the immediate upper binary level.

Let be two unsigned binary  numbers: $A = \{a_3, a_2, a_1, a_0\}$ and $B = \{b_3, b_2, b_1, b0\}$.
$$A + B = C$$
where: $C = \{c_3, c_2, c_1, c_0\}$ with $c_i$ computed by turn:\\
$c_0 = sum(a_0,b_0)$ and $carry0 = carry(a_0,b_0)$\\
$c_1 = sum(sum(a_1, b_1), carry0)$ and $carry1 = carry(a_1, b_1) \; | \; carry(sum(a_1, b_1), carry0)$\\
$c_2 = sum(sum(a_2, b_2), carry1)$ and $carry2 = carry(a_2, b_2) \; | \; carry(sum(a_2, b_2), carry1)$\\
$c_3 = sum(sum(a_3, b_3), carry2)$ and $carry3 = carry(a_3, b_3) \; | \; carry(sum(a_3, b_3), carry2)$\\
The operator $|$ is the bitwise OR.

Example, for two 8-bit unsigned integers:
\begin{verbatim}
       101 => 01100101 +
       142 => 10001110
              --------
       243 => 11110011
\end{verbatim}

\subsubsection{Subtracting binary numbers}

The subtract operation is reduced to the addition operation. It is performed by adding the negative number, i.e., $A - B = A + (-B)$.
Example, for two 8-bit signed integers:
\begin{verbatim}
        21 => 00010101 -
        34 => 00100010
              --------
              00010101 +
              11011101 +
              00000001
              --------
              11110010 +
              00000001
              --------
              11110011
\end{verbatim}
To check if the result is correct we must change the sign of the result:
\begin{verbatim}
              11110011 => 00001100 + 1 = 00001101 = 13 = -(21 - 34)
\end{verbatim}

\subsection{Hexadecimal Encoding}

\begin{verbatim}
     Decimal    Hexadecimal
        0:          0000
        1:          0001
        2:          0010
        3:          0011
        4:          0100
        5:          0101
        6:          0110
        7:          0111
        8:          1000
        9:          1001
        a/A:        1010
        b/B:        1011
        c/C:        1100
        d/D:        1101
        e/E:        1110
        f/F:        1111
\end{verbatim}

\begin{exemplu}
$$(2AC6)_{16} = (0010\_1010\_1100\_0110)_2 = (0010101011000110)_2 = 0010\_1010\_1100\_0110$$
$$(2AC6)_{16} = 2 \times 16^3 + (a)_{16} \times 16^2 + (c)_{16} \times 16^1 + 6 \times 16^0 = (10,950)_{10} = 10,950$$

Underscore are used only for easy read.

$\diamond$
\end{exemplu}

\subsection{Octal Encoding}

\begin{verbatim}
      Decimal      Octal
        0:          000
        1:          001
        2:          010
        3:          011
        4:          100
        5:          101
        6:          110
        7:          111
\end{verbatim}

\begin{exemplu}
$$(236)_8 = (010\_011\_110)_2 = (010011110)_2 =  010011110$$
$$(236)_8 = 2 \times 8^2 + 3 \times 8^1 + 6 \times 8^0 = (158)_{10} = 158$$

$\diamond$
\end{exemplu}

\subsection{Behavioral vs. Structural Descriptions}

The description of a digital circuit can be done in two ways:
\begin{itemize}
\item  by describing the functional behavior
\item  by describing the internal structure of the circuit
\end{itemize}
We exemplify the two methods in the case of the addition function.


\begin{exemplu}\label{adder}
Let us use the {\em System Verilog} HDL to describe an adder for 4-bit
numbers (see Figure \ref{first_ex}a). The description which
follows is a {\bf behavioral} one, because we know {\bf what} we
intend to design, but we do not know yet {\bf how} to design the
internal structure of an adder.

\placedrawing[h]{figures/01_04_4bit_adder.lp}{{\bf The first examples of  digital
systems.} {\small {\bf a.} The two 4-bit numbers adder, called
{\tt adder}. {\bf b.} The structure of an adder for 3 4-bit
numbers, called {\tt threeAdder}.} }{first_ex}

The Verilog code describing the module {\tt adder} is:
{\small
{\em
\begin{shaded*}
\begin{lstlisting}[language=Verilog, caption= , morekeywords={always_ff, always_comb, logic, loop, repeat, doInParallel, for}]
/*************************************************************************
File name:      adder.sv
Circuit name:   Adder
Description:    The module 'adder' has 2 4-bit inputs and one 4-bit output
                The circuit adds modulo 16; do not use or provide carry
                signal
*************************************************************************/
module adder(   output logic [3:0] out ,    // 4-bit output
                input  logic [3:0] in0 ,    // 4-bit input
                input  logic [3:0] in1 );   // 4-bit input

    assign out = in0 + in1;
endmodule
\end{lstlisting}
\end{shaded*}
}
}

The story just told by the previous Verilog module is: ``the 4-bit adder has two inputs, {\tt in0, in1}, one output, {\tt out}, and
its output is continuously assigned to the value obtained by
adding modulo 16 the two input numbers".

  $\diamond$
\end{exemplu}

What we just learned from the previous first simple example is
summarized in the following {\bf SystemVerilogSummary}.

\begin{verisum}{\em :

\hspace{-0.8cm} \rule[5mm]{16cm}{0.25mm} \vspace{-0.7cm}

\begin{description}
\item[/*]: begin comment
\item[*/]: end comment
\item[module]: keyword which indicates the beginning of the
description of a circuit as a module having the name which
immediately follows (in our example, the name is: {\tt adder})

\item[endmodule]: keyword which indicates the end of the module's
description which started with the previous keyword {\bf module}

\item[output]: keyword used to declare a terminal as an output (in
our example the terminal {\tt out} is declared as output)

\item[input]: keyword used to declare the terminal as an input (in
our example the terminals {\tt in0} and {\tt in1} are declared as
inputs)

\item[logic]: a data type with 4-state bits: {\tt 0, 1, x, z} where 0 means low, 1 means high, x means unknown, and z means an undriven net.

\item[initial]: a block which starts at the beginning of simulation

\item[begin ... end]: block delimiters

\item[\#]: indicate a raw value

\item[forever]: indicate an unending operation

\item[$\sim$]: negation operator

\item[assign]: keyword called the {\em continuous assignment},
used here to specify the function performed by the module (the
output {\tt out} takes continuously the value computed by adding
the two input numbers)

\item[(...)]: delimiters used to delimit the list of terminals
(external connections)

\item[,]: delimiter to separate each terminal within a list of
terminals

\item[;]: delimiter for end of line

\item[[ $ \dots$ ]]: delimiters which contains the definition of
the bits associated with a connection, for example [3:0] define
the number of bits for the three connections in the previous
example

\item[+]: the operator {\tt add}, the only one used in the
previous example.
\end{description}
\rule[5mm]{16cm}{0.25mm}
}
\end{verisum}

The description of a digital system is a {\bf hierarchical
construct} starting from a top module populated by modules, which
are similarly defined. The process continues until very simple
module are directly described. Thus, the functions $f$ and $g$ are specified by HDL programs
(in our case, in the previous example by a Verilog program).

The main characteristic of the digital design is {\bf modularity}. A
problem is decomposed in many simpler problems, which are solved
similarly, and so on until very simple problems are identified.
Modularity means also to define as many as possible identical
modules in each design. This allow to replicate many times the
same module, already designed and validated. {\em Many \& simple}
modules! Is the main slogan of the digital designer. Let's take
another example which uses as module the one just defined in the
previous example.

\begin{exemplu}
The previously exemplified module ({\tt adder}) will be used to
design a modulo 16 3-number adder, called {\tt threeAdder} (see Figure
\ref{first_ex}b). It adds 3 4-bit numbers providing a 4-bit result
(modulo 16 sum). Follows the {\bf structural} description:
{\small
{\em
\begin{shaded*}
\begin{lstlisting}[language=Verilog, caption= , morekeywords={always_ff, always_comb, logic, loop, repeat, doInParallel, for}]
/*************************************************************************
File name:      threeAdder.sv
Circuit name:   Three Input Adder
Description:    The module 'threeAdder' has 3 4-bit inputs and one 4-bit
                output. The circuit adds modulo 16 three numbers;
                do not provide carry output
*************************************************************************/
module threeAdder( output   logic [3:0]  out,
                   input    logic [3:0]  in0,
                   input    logic [3:0]  in1,
                   input    logic [3:0]  in2);

    logic [3:0] sum ;

    adder   inAdder(.out(sum),
                    .in0(in1),
                    .in1(in2)),
            outAdder(.out(out),
                     .in0(in0),
                     .in1(sum));
endmodule
\end{lstlisting}
\end{shaded*}
}
}

Two modules of {\tt adder} type (defined in the previous example)
are instantiated as {\tt inAdder}, {\tt outAdder}, they are
interconnected using the wire {\tt sum}, and are connected to the
terminals of the {\tt threeAdder} module. The resulting structure
computes the sum of three numbers. $\diamond$
\end{exemplu}

\begin{verisum}{\em :

\hspace{-0.8cm} \rule[5mm]{16cm}{0.25mm} \vspace{-0.7cm}

\begin{itemize}
\item How a previously defined module (in our example: {\tt
adder}) is two times instantiated using two different names ({\tt
inAdder} and {\tt outAdder} in our example)

\item A ``safe" way to allocate the terminals for a module
previously defined and instantiated inside the current module:
each original terminal name is preceded by a dot, and followed by
a parenthesis containing the name of the wire or of the terminal
where it is connected (in our example, {\tt outAdder( ...
.in1(sum))} means: the terminal {\tt in1} of the instance {\tt
outAdder} is connected to the wire {\tt sum})

\item The successive instantiations of the same module can be
separated by a ``,".
\end{itemize}
\rule[5mm]{16cm}{0.25mm} }
\end{verisum}

While the module {\tt adder} is a {\em behavioral} description,
the module {\tt threeAdder} is a {\em structural} one. The first
tells us {\em what} is the function of the module, and the second
tells us {\em how} its functionality is performed by using a
structure containing two instantiation of a previously defined
subsystems, and an internal connection.